\section{Why Algorithm Choice Determines Which Tail You Land In}
\label{sec:algorithm-choice}

\subsection{The B.0 Bakeoff}

The B.0 paper \citep{dellaLibera2026b0} compared three algorithmic approaches
to congressional redistricting: (1) minimum edge-cut (METIS), which is the
\ar{} algorithm; (2) adaptive bisection with compactness penalty; and
(3) a VRA-first multi-constraint approach that targets minority district
thresholds as a hard constraint.

The three algorithms produce similar plans in competitive states with moderate
sorting.  They diverge in sorted states and in states with concentrated
minority populations.

\subsection{Compactness Algorithms and the Republican Tail}

Algorithms that optimise compactness (including \ar{}) tend to land in
the right tail of the compactness distribution and the left tail (fewer
Democratic seats) of the partisan seat distribution in sorted states.
This is not a bug: it reflects the geographic reality that compact districts
in sorted states tend to have more Republican seats.

The mechanism: a compact algorithm draws boundaries that respect geographic
clustering.  Geographic clusters in sorted states are partisan clusters.
A compact algorithm will draw one or two large Democratic urban districts
and many smaller Republican rural/suburban districts, regardless of the
specific compactness objective used.

\subsection{Short-Burst Algorithms and Tail Risk}

Short-burst algorithms \citep{cannon2022voting} can produce extreme partisan
outcomes if the burst objective is correlated with partisanship.  For example,
a short-burst algorithm that maximises the number of competitive districts
(defined as those within 5 percentage points of 50\%) will tend to produce
plans that pack extreme partisans, which can produce outcomes at the
1st or 99th percentile of the partisan distribution.

This tail risk is a known limitation of short-burst methods.  It is absent
from \ar{} because the METIS objective has no partisan component.

\subsection{The Implication for Legal Arguments}

The algorithm choice is legally significant.  An algorithm with a
compactness objective will tend to produce Republican-leaning outcomes in
sorted states.  This does not mean it is biased---it means it is consistent.
The ensemble framework confirms that any \emph{compact} algorithm will
produce similar outcomes in sorted states; the bias, if any, belongs to
geography, not algorithm design.

The legal argument is therefore:
\begin{enumerate}
  \item The \ar{} algorithm has no partisan objective.
  \item Any compact algorithm produces similar outcomes in sorted states.
  \item The ensemble confirms that \ar{} outcomes are geographically typical
    for compact algorithms in all studied states.
  \item Therefore, the \ar{} plan is not a product of partisan intent.
\end{enumerate}

This argument is strengthened by the B.0 bakeoff, which shows that different
compact algorithms produce similar partisan outcomes, further confirming that
the outcomes are geographic rather than algorithmic.
